%%%%%%%%%%%%%%%%%%%%%%%%%%%%%%%%%%%%%%%%%%%%%%%%%%%%%%%%%%%%%%%%%%%%%%%%%%%%%%%%%%%
%% Apresentação do TCC - Proposta de Reforma Previdenciária                      %%
%% author: Julio Cesar Ferreira Lima                                             %%
%% Universidade Federal do Ceará (UFC)                                           %%
%%%%%%%%%%%%%%%%%%%%%%%%%%%%%%%%%%%%%%%%%%%%%%%%%%%%%%%%%%%%%%%%%%%%%%%%%%%%%%%%%%%
\documentclass{libs/ufc_format}
% Inserting the preamble file with the packages
\input{libs/preamble.tex}
% Inserting the references file
\bibliography{references.bib}

% Title
\title[Reforma Previdenciária via Blockchain]{\huge\textbf{Proposta de Reforma Previdenciária}}
% Subtitle
\subtitle{Um Mercado Financeiro Poupador via Blockchain como Alternativa ao INSS e FGTS}
% Author of the presentation
\author{Julio Cesar Ferreira Lima}
% Institute's Name
\institute[UFC]{
    % email for contact
    \normalsize{\email{julio_flima@hotmail.com}}
    \newline
    % Department Name
    \department{Departamento de Engenharia Elétrica}
    \newline
    % university name
    \ufc
}
% date of the presentation
\date{2026}


%%%%%%%%%%%%%%%%%%%%%%%%%%%%%%%%%%%%%%%%%%%%%%%%%%%%%%%%%%%%%%%%%%%%%%%%%%%%%%%%%%
%% Start Document of the Presentation                                           %%               
%%%%%%%%%%%%%%%%%%%%%%%%%%%%%%%%%%%%%%%%%%%%%%%%%%%%%%%%%%%%%%%%%%%%%%%%%%%%%%%%%%
\begin{document}
% insert the code style
\input{libs/code_style}

%% ---------------------------------------------------------------------------
% First frame (with tile, subtitle, ...)
\begin{frame}{}
    \maketitle
\end{frame}

%% ---------------------------------------------------------------------------
% Second frame
\begin{frame}{Sumário}
    \begin{multicols}{2}
        \tableofcontents
    \end{multicols}
\end{frame}

%% ---------------------------------------------------------------------------
%% INTRODUÇÃO
%% ---------------------------------------------------------------------------
\section{Introdução}
\begin{frame}{Contextualização}
    \begin{itemize}
        \item O Brasil enfrenta \textbf{desafios econômicos estruturais} que limitam seu desenvolvimento há décadas
        \item O foco excessivo em corrupção e polarização desvia a atenção dos problemas verdadeiramente estruturantes:
        \begin{itemize}
            \item Infraestrutura deficiente
            \item Baixa taxa de poupança nacional
            \item Crédito historicamente caro
        \end{itemize}
    \end{itemize}
    
    \vspace{0.3cm}
    
    \begin{alertblock}{Problema Central}
        Os gastos com previdência e assistência social representam aproximadamente \textbf{54\% das despesas primárias} do governo federal.
    \end{alertblock}
\end{frame}

%% ---------------------------------------------------------------------------
\begin{frame}{O Sistema Previdenciário Atual}
    \begin{columns}
        \begin{column}{0.5\textwidth}
            \begin{block}{INSS}
                \begin{itemize}
                    \item Regime de repartição simples
                    \item Déficit de R\$ 318,4 bi (2023)
                    \item Trabalhadores ativos financiam aposentados
                \end{itemize}
            \end{block}
        \end{column}
        \begin{column}{0.5\textwidth}
            \begin{block}{FGTS}
                \begin{itemize}
                    \item 8\% do salário depositado
                    \item Rentabilidade inferior à inflação
                    \item Perda real de até 40\% em 30 anos
                \end{itemize}
            \end{block}
        \end{column}
    \end{columns}
    
    \vspace{0.4cm}
    
    \simplebox{O trabalhador \textbf{não possui propriedade real} sobre suas contribuições!}
\end{frame}

%% ---------------------------------------------------------------------------
\subsection{Justificativa e Objetivos}
\begin{frame}{Justificativa}
    \begin{itemize}
        \item Necessidade urgente de reformar o sistema previdenciário brasileiro
        \item Reformas implementadas até o momento têm sido \textbf{paliativas}
        \item A tecnologia blockchain oferece oportunidade única de implementar um sistema:
        \begin{itemize}
            \item Transparente
            \item Imutável
            \item Descentralizado
        \end{itemize}
    \end{itemize}
    
    \vspace{0.3cm}
    
    \successbox{Garantir a propriedade individual dos recursos previdenciários com mecanismos de controle contra uso inadequado.}
\end{frame}

%% ---------------------------------------------------------------------------
\begin{frame}{Objetivos}
    \textbf{Objetivo Geral:}
    \begin{itemize}
        \item Propor um modelo de mercado financeiro poupador baseado em blockchain como alternativa ao INSS e FGTS
    \end{itemize}
    
    \vspace{0.3cm}
    
    \textbf{Objetivos Específicos:}
    \begin{enumerate}
        \item Analisar os problemas estruturais do sistema previdenciário brasileiro atual
        \item Estudar o modelo australiano de \textit{Superannuation}
        \item Propor uma arquitetura de sistema blockchain para gestão previdenciária
        \item Identificar benefícios e desafios da implementação
    \end{enumerate}
\end{frame}

%% ---------------------------------------------------------------------------
%% FUNDAMENTAÇÃO TEÓRICA
%% ---------------------------------------------------------------------------
\section{Fundamentação Teórica}

\subsection{Modelo Australiano}
\begin{frame}{Modelo Australiano de Superannuation}
    \begin{columns}
        \begin{column}{0.6\textwidth}
            \begin{itemize}
                \item Sistema de capitalização individual obrigatória desde 1992
                \item Empregadores contribuem com \textbf{11\% do salário}
                \item Trabalhador escolhe fundo e estratégia de investimento
                \item Acesso apenas na idade de preservação (55-60 anos)
            \end{itemize}
        \end{column}
        \begin{column}{0.4\textwidth}
            \begin{exampleblock}{Resultados}
                \begin{itemize}
                    \item \textbf{3,5 trilhões AUD} acumulados
                    \item \textbf{170\% do PIB} australiano
                \end{itemize}
            \end{exampleblock}
        \end{column}
    \end{columns}
    
    \vspace{0.3cm}
    
    \successbox{Demonstra que sistemas de capitalização individual \textbf{podem ser implementados com sucesso}.}
\end{frame}

%% ---------------------------------------------------------------------------
\subsection{Tecnologia Blockchain}
\begin{frame}{Fundamentos de Blockchain}
    \begin{block}{Características Principais}
        \begin{itemize}
            \item \textbf{Descentralização}: Dados replicados em múltiplos nós
            \item \textbf{Imutabilidade}: Transações não podem ser alteradas
            \item \textbf{Transparência}: Regras públicas e verificáveis
            \item \textbf{Segurança criptográfica}: Proteção por algoritmos robustos
        \end{itemize}
    \end{block}
    
    \vspace{0.3cm}
    
    \begin{alertblock}{Inovação Central}
        Resolve o \textbf{Problema dos Generais Bizantinos}: múltiplos participantes alcançam consenso sem autoridade central.
    \end{alertblock}
\end{frame}

%% ---------------------------------------------------------------------------
\begin{frame}{Smart Contracts e Multi-sig}
    \begin{columns}
        \begin{column}{0.5\textwidth}
            \begin{block}{Smart Contracts}
                \begin{itemize}
                    \item Programas autoexecutáveis
                    \item Executam quando condições são satisfeitas
                    \item Automação de regras complexas
                \end{itemize}
            \end{block}
        \end{column}
        \begin{column}{0.5\textwidth}
            \begin{block}{Carteiras Multi-sig}
                \begin{itemize}
                    \item Múltiplas chaves para autorizar transação
                    \item Exemplo: 2-de-3 chaves
                    \item Controle compartilhado
                \end{itemize}
            \end{block}
        \end{column}
    \end{columns}
    
    \vspace{0.4cm}
    
    \simplebox{Multi-sig permite implementar \textbf{governança compartilhada} onde nenhuma parte possui controle total.}
\end{frame}

%% ---------------------------------------------------------------------------
%% METODOLOGIA
%% ---------------------------------------------------------------------------
\section{Metodologia}
\begin{frame}{Classificação da Pesquisa}
    \begin{itemize}
        \item \textbf{Natureza}: Aplicada (solução para problema específico)
        \item \textbf{Objetivos}: Exploratória e propositiva
        \item \textbf{Abordagem}: Qualitativa
    \end{itemize}
    
    \vspace{0.3cm}
    
    \begin{block}{Procedimentos Técnicos}
        \begin{enumerate}
            \item \textbf{Pesquisa Bibliográfica}: Bases acadêmicas e documentação técnica
            \item \textbf{Pesquisa Documental}: Legislação, relatórios oficiais
            \item \textbf{Análise Comparativa}: Brasil vs. Austrália
        \end{enumerate}
    \end{block}
\end{frame}

%% ---------------------------------------------------------------------------
\begin{frame}{Etapas de Desenvolvimento}
    \begin{enumerate}
        \item \textbf{Diagnóstico do Sistema Atual}
        \begin{itemize}
            \item Estrutura de contribuições e benefícios
            \item Evolução histórica do déficit
        \end{itemize}
        
        \item \textbf{Estudo de Modelos Internacionais}
        \begin{itemize}
            \item Modelo australiano de Superannuation
        \end{itemize}
        
        \item \textbf{Estudo da Tecnologia Blockchain}
        \begin{itemize}
            \item Arquitetura de redes distribuídas
            \item Smart contracts e multi-sig
        \end{itemize}
        
        \item \textbf{Elaboração da Proposta}
        
        \item \textbf{Validação Conceitual}
    \end{enumerate}
\end{frame}

%% ---------------------------------------------------------------------------
%% PROPOSTA DO SISTEMA
%% ---------------------------------------------------------------------------
\section{Proposta do Sistema}

\subsection{Visão Geral}
\begin{frame}{Princípios Fundamentais}
    \begin{exampleblock}{Visão Geral}
        Mercado financeiro poupador que substitui gradualmente INSS e FGTS, permitindo que trabalhadores invistam em empresas brasileiras.
    \end{exampleblock}
    
    \vspace{0.3cm}
    
    \begin{itemize}
        \item \textbf{Propriedade individual}: Recursos pertencem ao trabalhador
        \item \textbf{Investimento produtivo}: Recursos direcionados a empresas reais
        \item \textbf{Governança compartilhada}: Controle dividido (trabalhador + governo)
        \item \textbf{Transparência nas regras}: Código auditável
        \item \textbf{Automatização}: Regras executadas por smart contracts
    \end{itemize}
\end{frame}

%% ---------------------------------------------------------------------------
\subsection{Arquitetura}
\begin{frame}{Arquitetura do Sistema}
    \begin{figure}
        \centering
        \begin{tikzpicture}[scale=0.8,
            node distance=1.2cm,
            box/.style={rectangle, draw, minimum width=3.5cm, minimum height=0.8cm, align=center, font=\small},
            arrow/.style={->, >=stealth, thick}
        ]
            % Camadas
            \node[box, fill=blue!20] (user) {Camada de Usuário\\(Interface Web/Mobile)};
            \node[box, fill=green!20, below=of user] (app) {Camada de Aplicação\\(Smart Contracts)};
            \node[box, fill=orange!20, below=of app] (chain) {Camada de Blockchain\\(Rede Distribuída)};
            
            % Conexões
            \draw[arrow] (user) -- (app);
            \draw[arrow] (app) -- (chain);
            
            % Atores
            \node[box, fill=gray!20, right=1.5cm of user] (trab) {Trabalhador};
            \node[box, fill=gray!20, right=1.5cm of app] (gov) {Governo};
            \node[box, fill=gray!20, right=1.5cm of chain] (emp) {Empresas};
            
            \draw[arrow] (trab) -- (user);
            \draw[arrow] (gov) -- (app);
            \draw[arrow] (emp) -- (chain);
        \end{tikzpicture}
    \end{figure}
    
    \simplebox{Blockchain \textbf{permissionada} operada por Banco Central e instituições reguladas.}
\end{frame}

%% ---------------------------------------------------------------------------
\subsection{Governança Multi-sig}
\begin{frame}{Sistema de Governança Multi-sig (2-de-2)}
    \begin{columns}
        \begin{column}{0.5\textwidth}
            \begin{block}{Chaves}
                \begin{itemize}
                    \item \textbf{Chave do Trabalhador}: Dispositivo pessoal seguro
                    \item \textbf{Chave do Governo}: Órgão regulador (ex: BC)
                \end{itemize}
            \end{block}
        \end{column}
        \begin{column}{0.5\textwidth}
            \begin{alertblock}{Garantias}
                \begin{itemize}
                    \item Governo \textbf{não pode} mover recursos sozinho
                    \item Trabalhador \textbf{não pode} violar regras
                \end{itemize}
            \end{alertblock}
        \end{column}
    \end{columns}
    
    \vspace{0.4cm}
    
    \begin{block}{Fluxo de Autorização}
        \begin{enumerate}
            \item Trabalhador inicia solicitação (assina com sua chave)
            \item Smart contract verifica regras
            \item Sistema governamental co-assina automaticamente se válido
            \item Transação executada e registrada na blockchain
        \end{enumerate}
    \end{block}
\end{frame}

%% ---------------------------------------------------------------------------
\subsection{Regras Operacionais}
\begin{frame}{Regras Operacionais}
    \begin{table}
        \centering
        \small
        \begin{tabular}{p{3.5cm}p{6cm}}
            \hline
            \textbf{Regra} & \textbf{Descrição} \\
            \hline
            Operações mensais & Máximo 1 lote de compra/venda por mês \\
            Ativos elegíveis & Apenas empresas/ativos brasileiros \\
            Limites de concentração & Configuráveis na adesão, imutáveis após aceite \\
            Saque & Apenas na aposentadoria ou exceções legais \\
            Governança & Multi-sig 2-de-2 (trabalhador + governo) \\
            \hline
        \end{tabular}
    \end{table}
    
    \vspace{0.3cm}
    
    \successbox{Restrições evitam comportamento especulativo e garantem uso previdenciário.}
\end{frame}

%% ---------------------------------------------------------------------------
\begin{frame}{Universo de Investimentos}
    \begin{block}{Ativos Elegíveis (apenas brasileiros)}
        \begin{itemize}
            \item Ações de empresas listadas na B3
            \item Debêntures de empresas brasileiras
            \item Fundos de investimento em infraestrutura (FI-Infra)
            \item Fundos imobiliários (FIIs)
            \item Títulos públicos federais
        \end{itemize}
    \end{block}
    
    \vspace{0.3cm}
    
    \begin{exampleblock}{Objetivos da Restrição}
        \begin{itemize}
            \item Fomentar o mercado de capitais brasileiro
            \item Financiar o crescimento de empresas nacionais
            \item Evitar evasão de divisas
            \item Criar ciclo virtuoso de poupança interna
        \end{itemize}
    \end{exampleblock}
\end{frame}

%% ---------------------------------------------------------------------------
\subsection{Herança}
\begin{frame}{Mecanismo de Herança}
    \begin{alertblock}{Diferencial Fundamental}
        Em caso de falecimento, \textbf{100\% do patrimônio} é transferido para herdeiros designados --- diferente do INSS atual!
    \end{alertblock}
    
    \vspace{0.3cm}
    
    \begin{columns}
        \begin{column}{0.5\textwidth}
            \begin{block}{Cadastro de Beneficiários}
                \begin{itemize}
                    \item Endereço da carteira do sucessor
                    \item Percentual destinado a cada um
                    \item \textbf{Autonomia absoluta} do trabalhador
                \end{itemize}
            \end{block}
        \end{column}
        \begin{column}{0.5\textwidth}
            \begin{block}{Processo de Transferência}
                \begin{enumerate}
                    \item Família apresenta certidão + chave
                    \item Patrimônio transferido
                    \item Recursos bloqueados até aposentadoria do herdeiro
                \end{enumerate}
            \end{block}
        \end{column}
    \end{columns}
\end{frame}

%% ---------------------------------------------------------------------------
%% BENEFÍCIOS E DESAFIOS
%% ---------------------------------------------------------------------------
\section{Benefícios e Desafios}

\begin{frame}{Benefícios Esperados}
    \begin{columns}
        \begin{column}{0.33\textwidth}
            \begin{block}{Para o Trabalhador}
                \begin{itemize}
                    \item Propriedade real
                    \item Rentabilidade superior
                    \item Herança garantida
                    \item Transparência
                \end{itemize}
            \end{block}
        \end{column}
        \begin{column}{0.33\textwidth}
            \begin{block}{Para a Economia}
                \begin{itemize}
                    \item Maior poupança
                    \item Mercado de capitais
                    \item Menor custo de capital
                    \item Investimentos produtivos
                \end{itemize}
            \end{block}
        \end{column}
        \begin{column}{0.33\textwidth}
            \begin{block}{Para o Estado}
                \begin{itemize}
                    \item Fim do déficit
                    \item Recursos liberados
                    \item Menor carga tributária
                    \item Modernização
                \end{itemize}
            \end{block}
        \end{column}
    \end{columns}
\end{frame}

%% ---------------------------------------------------------------------------
\begin{frame}{Desafios e Riscos}
    \begin{columns}
        \begin{column}{0.5\textwidth}
            \begin{alertblock}{Desafios Técnicos}
                \begin{itemize}
                    \item Escalabilidade da blockchain
                    \item Segurança cibernética
                    \item Integração com sistemas legados
                    \item Educação financeira
                \end{itemize}
            \end{alertblock}
        \end{column}
        \begin{column}{0.5\textwidth}
            \begin{alertblock}{Desafios Políticos}
                \begin{itemize}
                    \item Resistência de grupos beneficiados
                    \item Reforma constitucional
                    \item Coordenação entre órgãos
                    \item Comunicação com população
                \end{itemize}
            \end{alertblock}
        \end{column}
    \end{columns}
    
    \vspace{0.3cm}
    
    \simplebox{Riscos de mercado: volatilidade e perdas em crises --- necessários mecanismos de proteção.}
\end{frame}

%% ---------------------------------------------------------------------------
%% CONCLUSÃO
%% ---------------------------------------------------------------------------
\section{Conclusão}
\begin{frame}{Conclusões}
    \begin{itemize}
        \item O sistema previdenciário brasileiro apresenta \textbf{problemas estruturais significativos}
        \item O modelo australiano demonstrou \textbf{viabilidade} de capitalização individual
        \item A tecnologia blockchain oferece \textbf{mecanismos técnicos} para transparência e governança
    \end{itemize}
    
    \vspace{0.3cm}
    
    \begin{exampleblock}{Sistema Proposto}
        \begin{itemize}
            \item Trabalhadores são \textbf{proprietários reais} de seus recursos
            \item Recursos investidos em \textbf{empresas brasileiras}
            \item Governança \textbf{compartilhada} via multi-sig 2-de-2
            \item Patrimônio \textbf{integralmente herdável}
        \end{itemize}
    \end{exampleblock}
\end{frame}

%% ---------------------------------------------------------------------------
\begin{frame}{Trabalhos Futuros}
    \begin{columns}
        \begin{column}{0.5\textwidth}
            \begin{block}{Desenvolvimento Técnico}
                \begin{itemize}
                    \item Protótipo funcional (testnet)
                    \item Análise de escalabilidade
                    \item Interfaces acessíveis
                \end{itemize}
            \end{block}
            
            \vspace{0.3cm}
            
            \begin{block}{Análise Quantitativa}
                \begin{itemize}
                    \item Simulação de impacto fiscal
                    \item Modelagem atuarial
                \end{itemize}
            \end{block}
        \end{column}
        \begin{column}{0.5\textwidth}
            \begin{block}{Aspectos Jurídicos}
                \begin{itemize}
                    \item Alterações constitucionais
                    \item Marco regulatório
                \end{itemize}
            \end{block}
            
            \vspace{0.3cm}
            
            \begin{block}{Aspectos Sociais}
                \begin{itemize}
                    \item Pesquisa de aceitação
                    \item Programas de educação financeira
                \end{itemize}
            \end{block}
        \end{column}
    \end{columns}
\end{frame}

%% ---------------------------------------------------------------------------
% Reference frames
\begin{frame}[allowframebreaks]
    \frametitle{Referências}
    \printbibliography
\end{frame}

%% ---------------------------------------------------------------------------
% Final frame
\begin{frame}{}
    \centering
    \huge{\textbf{\example{Obrigado pela Atenção!}}}
    
    \vspace{1cm}
    
    \normalsize{Julio Cesar Ferreira Lima}
    
    \small{julio\_flima@hotmail.com}
    
    \vspace{0.5cm}
    
    \small{Orientador: Prof. Dr. José Cláudio do Nascimento}
\end{frame}

\end{document}